
\section{动机例：胀开平面原点}

我们将推广下述概念，它将相应于仿射平面 $\mathbb{A}^2_k$ 原点的「\trBlowup」（\mbox{练习~10.3.F}）。我们的讨论将是非正式的。考虑 $\mathbb{A}^2\times\mathbb{P}^1$ 中相应于如下条件的子集。我们把 $\mathbb{P}^1$ 解读为过原点诸直线的参数化空间。考虑\trSubvariety
\[
\mathrm{Bl}_{(0,0)}\mathbb{A}^2:=\{(p\in\mathbb{A}^2,[\ell]\in\mathbb{P}^1)\;:\;p\in\ell\},
\]
它所给出的资料是：平面上的一个点 $p$，与一条同时包含 $p$ 和原点的直线 $\ell$。从代数上看：取 $x$、$y$ 为 $\mathbb{A}^2$ 上的\trCoordinate，$X$、$Y$ 为 $\mathbb{P}^1$ 上的射影\trCoordinate（与 $x$、$y$「相应」）；我们将考虑 $\mathbb{A}^2\times\mathbb{P}^1$ 中相应于 $xY-yX=0$ 的子集 $\mathrm{Bl}_{(0,0)}\mathbb{A}^2$。我们有一张有用的图：
\[
\begin{tikzcd}
\mathrm{Bl}_{(0,0)}\mathbb{A}^2\ar[r,hook]\ar[rd,"\beta"'] & \mathbb{A}^2\times\mathbb{P}^1\ar[r]\ar[d] & \mathbb{P}^1 \\
& \mathbb{A}^2 &
\end{tikzcd}
\]

你可以直接验证 $\mathrm{Bl}_{(0,0)}\mathbb{A}^2$ 在 $k$ 上\trSmooth{}（\mbox{定义~13.2.4}或\mbox{21.2.28}）（你现在可以让\mbox{练习~10.3.F}之后的那段话变得精确），但这里给出一个非正式的论证，利用\trProjection{} $\mathrm{Bl}_{(0,0)}\mathbb{A}^2\to\mathbb{P}^1$。\trProjectiveLine{} $\mathbb{P}^1$ 是\trSmooth{}的，而对于 $\mathbb{P}^1$ 中的每个点 $[\ell]$，直线 $\ell$ 上的诸点都有一个光滑的选取。这样，我们便是在通过 $\mathbb{P}^1$ 上的一个\trFibration{}来验证\trSmoothness。

接下来我们考虑投向 $\mathbb{A}^2$ 的\trProjection{} $\beta\colon\mathrm{Bl}_{(0,0)}\mathbb{A}^2\to\mathbb{A}^2$。它在原点之外是一个\trIsomorphism。粗略地说，若 $p$ 不是原点，则恰有一条直线同时包含 $p$ 和原点。另一方面，若 $p$ 就是原点，则同时包含 $p$ 和原点的直线有整整一个 $\mathbb{P}^1$ 那么多。于是 $(0,0)$ 的原像是一条\trCurve，从而是一个\trDivisor（一个\trEffectiveCartierDivisor，因为胀开所得的曲面是\trRegular{}的）。它称作这次\trBlowup{}的\trExceptionalDivisor。

倘若我们有一条在原点处\trSingular{}的\trCurve{} $C\subset\mathbb{A}^2$，那么借助这次\trBlowup，取 $C\setminus\{(0,0)\}$ 在 $\mathrm{Bl}_{(0,0)}\mathbb{A}^2$ 中的\trClosure，便有望把它部分地\trDesingularization。（一个\trVariety{} $X$ 的\textbf{\trDesingularization}（desingularization）或\textbf{\trResolutionOfSingularities}（resolution of singularities），是一个形如 $\widetilde{X}\to X$ 的\trProper{}\trBirational{}\trMorphism，其中 $\widetilde{X}$ 是\trRegular{}概形。）例如，考虑曲线 $y^2=x^3+x^2$，它除去原点处的一个\trNode{}之外处处\trRegular。我们可以取这条曲线去掉原点之后的原像，再取这一轨迹在\trBlowup{}中的闭包，便会得到一条正则曲线；结点在「楼下」的两个分支，到「楼上」就被分开了。（等到我们把各种概念恰当地定义好之后，你可以在\mbox{练习~22.4.B}中核验这件事。所得结果将称为这条曲线的\emph{\trProperTransform}（或称\emph{strict transform}，严格变换）。）我们出于许多理由关心\trDesingularization。由于我们对正则\trCurve{}的理解已经相当透彻，我们可望通过一条\trCurve{}的\trDesingularization{}来理解其他的曲线。这一哲学在高维情形同样成立。

更一般地，我们可以把 $\mathbb{A}^n$ 在原点处胀开（或者更不正式地说，「胀开原点」），得到 $\mathbb{A}^n\times\mathbb{P}^{n-1}$ 的一个\trSubvariety。从代数上看，若 $x_1,\dots,x_n$ 是 $\mathbb{A}^n$ 上的\trCoordinate，$X_1,\dots,X_n$ 是 $\mathbb{P}^{n-1}$ 上的射影\trCoordinate，则\trBlowup{} $\mathrm{Bl}_{\vec 0}\mathbb{A}^n$ 由诸方程 $x_iX_j-x_jX_i=0$ 给出。同样地，它是\trSmooth{}的：$\mathbb{P}^{n-1}$ 是\trSmooth{}的，而对于每个点 $[\ell]\in\mathbb{P}^{n-1}$，$p\in\ell$ 都有光滑的选取。

我们还可以更进一步：在上一个例子中添上 $m$ 个新变元 $x_{n+1},\dots,x_{n+m}$，从而把 $\mathbb{A}^{n+m}$ 沿一个坐标 $m$ 维平面 $\mathbb{A}^m$ 胀开；我们便得到 $\mathbb{A}^{n+m}\times\mathbb{P}^{n-1}$ 的一个子集。

由于在复几何中，\trManifold{}的子流形局部地「看起来像」$n$ 维空间中的坐标 $m$ 维平面，你或许会想象，我们可以把它推广为胀开正则\trVariety{}的正则\trSubvariety。在把这说精确的过程中，我们将意外地把这个概念大大推广，定义概形中任意\trFiniteType{}\trQC{}\trIdealSheaf{}的\trBlowup。一般而言，\trBlowup{}未必有「在正则之物内部胀开正则之物」那一情形下那样直观的描述——它可能对概形大施暴力——但即便如此，它仍然非常有用。

我们的描述将只依赖于被胀开的那个\trClosedSubscheme，而不依赖于坐标。这弥补了第一个例子——在原点处胀开平面——中本已存在的一个缺陷：倘若我们为平面选取不同的坐标（保持原点作为一个\trClosedSubscheme{}不变），我们并不会因此得到两个互不相同的胀开，这一点绝非显然。

一如常例，理解\trBlowup{}概念有两种方式，各有其适用的场合。第一种是通过\trUniversalProperty，它让你不费吹灰之力便能证出某些事情。第二种是显式构造，它让你亲手弄脏双手、把东西算出来（并且，比如说，蕴含了\trBlowup{}\trMorphism{}是\trProjective{}的这一事实）。

这里的动机例子也许看似一个非常特殊的情形，但只要你对 $n$ 维空间中原点的\trBlowup{}理解得足够透彻，你就会理解一般的\trBlowup。
